\chapter{Sistemas dinámicos}

\section{Definición de sistemas}

\begin{figure}[h]
\centering
\includegraphics[scale=0.7]{sistema.png}
\end{figure}
Los sistemas son modelos matemáticos (ecuaciónes) que relacionan variables de entrada con variables de salida. Pueden ser representados gráficamente por diagramas de bloques. \\
Se utiliza matemática para ``modelar'' sistemas por las siguientes razones:
\begin{itemize}
    \item[(1)] Análisis
    \item[(2)] Diseño
    \item[(3)] Control y monitoreo
\end{itemize}


\section{Tipos de sistemas}
\subsection{Estáticos vs Dinámicos}
\begin{itemize}
    \item Sistema Estático: respuesta solo depende de la excitación actual (independiente del tiempo).
    \begin{figure}[h]
    \centering
    \includegraphics[scale=0.7]{estatico.png}
    \end{figure}
    
    \item Sistema Dinámico: respuesta actual depende de la excitación actual y de la ``historia'' de la respuesta.
    \begin{figure}[h]
    \centering
    \includegraphics[scale=0.7]{dinamico.png}
    \end{figure}
    \par Si la respuesta del sistema dinámico está basada en la excitación actual y pasada, entonces el sistema se conoce como \textcolor{red}{causal} (todos los sistemas físicos).\\
    \par Si la respuesta del sistema dinámico está basada en excitación actual, pasada y \textcolor{ForestGreen}{futura}, entonces el sistema es \textcolor{red}{anticausal} (teóricos).
\end{itemize}

\subsection{SISO vs Multivariable}
\begin{itemize}
    \item SISO: Single Input, Single Output.
    \begin{figure}[h]
    \centering
    \includegraphics[scale=0.7]{SISO.png}
    \end{figure}
    \begin{itemize}
        \item $u$: señal escalar
        \item $y$: señal escalar
        \item $f(u)=y$: función escalar
    \end{itemize}
    
    \item SIMO: Single Input, Multi Output.
    \begin{figure}[h]
    \centering
    \includegraphics[scale=0.7]{SIMO.png}
    \end{figure}
    \begin{itemize}
        \item $u$: señal escalar
        \item $\vec{y}$: señal vectorial
    \end{itemize}
    
    \item MISO: Multi Input, Simple Output.
    \begin{figure}[h]
    \centering
    \includegraphics[scale=0.7]{MISO.png}
    \end{figure}
    \begin{itemize}
        \item $\vec{u}$: señal vectorial
        \item $y$: señal escalar
    \end{itemize}
    
    \item MIMO: Multi Input, Multi Output.
    \begin{figure}[h]
    \centering
    \includegraphics[scale=0.7]{MIMO.png}
    \end{figure}
    \begin{itemize}
        \item $\vec{u}$: señal vectorial
        \item $\vec{y}$: señal vectorial
        \item $f(\vec{u})=\vec{y}$: función vectorial
    \end{itemize}
\end{itemize}

\subsection{Lineal vs No lineal}
\begin{itemize}
    \item Un \textcolor{red}{sistema lineal} es tal que cumple con el principio de \textcolor{red}{superposición}.
    \begin{figure}[h]
    \centering
    \includegraphics[scale=0.7]{sistemayu.png}
    \end{figure}
    \begin{itemize}
        \item Aditividad: Si $u_1(t) \rightarrow y_1(t)$, $u_2(t) \rightarrow y_2(t)$, entonces: $u(t) = u_1(t) + u_2(t) \rightarrow y(t) = y_1(t) + y_2(t)$
        \item Homogeneidad/Proporcionalidad: Si $u(t) \rightarrow y(t)$, entoncecs $\textcolor{red}{\alpha}u(t) \rightarrow \textcolor{red}{\alpha}y(t)$.
    \end{itemize}
    \item De no cumplir con las propiedades del sistema lineal, entonces el sistema es \textcolor{red}{no lineal}.
\end{itemize}
\par \underline{Ejemplo}: resorte no lineal.\\
\begin{figure}[h]
    \centering
    \includegraphics[scale=0.7]{ej_nolineal.png}
    \end{figure}
Se tiene que: 
\begin{align*}
f_s &= ku^2
p_1 &\rightarrow u_1\\  
p_2 &\rightarrow u_2\\
\end{align*}
La pregunta es, ¿Se cumple la condición de aditividad?
\begin{align*}
p_1+p_2 &\xrightarrow{?} u_1+u_2
\end{align*}
La respuesta es \textcolor{red}{no}, ya que:
\begin{align*}
    f_s = k(u_1+u_2)^2 \\
    p_1 + p_2 & = ku_1^2+ku_2^2+2ku_1u_2\\
    p_1 + p_2 \neq p_1 + p_2 + \textcolor{red}{2ku_1u_2}
\end{align*}

\subsection{Invariante tiempo vs variante tiempo}
\begin{itemize}
    \item Sistema invariantes en el tiempo $\longrightarrow$ parámetros del modelo son constantes en el tiempo.\\
    \underline{Ejemplo}: Sistemas de 1GDL con masa, rigidez y amortiguamiento constante.
    \begin{figure}[h]
    \centering
    \includegraphics[scale=0.7]{invariante.png}
    \end{figure}
    
    \item Sistema variante en el tiempo $\longrightarrow$ parámetros del modelo varían en el tiempo.\\
    \underline{Ejemplo}: Puentes con tráfico (la masa varía).
\end{itemize}

\subsection{Determinístico vs Estocástico}
Si el input del sistema es conocido (y el sistema no incorpora incertidumbre), entonces el output será predecible $\longrightarrow$ \textcolor{red}{Determinístico} \\ 
\underline{Ejemplo}: Simular múltiples veces el mismo registro (input) sobre el mismo edifico (sistema), se obtienen siempre las mismas respuestas (output). \\

Si el input del sistema posee incertidumbre (y/o el sistema incorpora incertidumbre), entonces el output también tiene incertidumbre y no es predecible $\longrightarrow$ \textcolor{red}{Estocástico} \\
\underline{Ejemplo}: Simular peatones diferentes (input con incertidumbre en masa del peatón, velocidad al caminar y frecuencia de pasos) sobre un mismo puente (sistema), se obtendrán respuestas diferentes (outputs)\\

% María: Si los parámetros del modelo son conocidos $\longrightarrow$ \textcolor{red}{Determinístico}.\\
% María: Pero, si existe incertidumbre en alguno de los parámetros, o si no se puede predecir de manera precisa la respuesta del sistema $\longrightarrow$ \textcolor{red}{Estocástico}.

\subsection{Continuo vs Discreto}
\begin{itemize}
    \item Sistema Continuo: entrada y salida son función del tiempo.
    \begin{figure}[h]
    \centering
    \includegraphics[scale=0.7]{continuo.png}
    \end{figure}
    \item Sistema Discreto: señales han sido discretizadas en el tiempo.\\
    % Ale: Reemplazar el ejemplo por siguiente línea:
    % Ejemplo: Registro de aceleraciones?
    \underline{Ejemplo}: Interés de una cuenta bancaria.
    \begin{align*}
        y(k+1)=(1+r)y(k) \longrightarrow \text{Ecuación de diferencias finitas}
    \end{align*}
    \begin{itemize}
        \item $y(k+1)$: balance cuenta año $k+1$
        \item $r$: interés compuesto anual
        \item $y(k)$: balance cuenta año $k$
    \end{itemize}
\end{itemize}

\section{``Time sampling''} % Reemplazar por MUESTREO?
Las señales de los sistemas físicos son continuas, pero que al momento de ser "observadas" por dispositivos digitales, estas señales se discretizan en función del tiempo.
\begin{figure}[h]
    \centering
    \includegraphics[scale=0.7]{time_sampling.png}
    \end{figure}



\section{Respuesta de sistemas LTI en tiempo continuo}
LTI: Linear, Time-Invariant. \\

\begin{figure}[h] % Borrar ?, ya está arriba
\centering
\includegraphics[scale=0.7]{sistemayu.png}
\end{figure}

Los sistemas LTI, son aquellos lineales cuyas propiedades y parámetros no varían en el tiempo, es decir, que los coeficientes de las ecuaciones diferenciales de los modelos matemáticos de los sistemas dinámicos son constantes en el tiempo $a_i=cte$.
\begin{align*}
    \frac{d^ny}{dt^n}+a_{n-1}\frac{d^{n-1}y}{dt^{n-1}}+...+a_2\frac{d^2y}{dt^2}+a_1\frac{dy}{dt}+a_0y=u(t)
\end{align*}
\begin{itemize}
    \item $u(t)$: excitación
    \item $y(t)$: respuesta
\end{itemize}
Dado un $u(t)$, la respuesta se obtiene al resolver la ecuación diferencial:
\begin{align*}
    y(t)=y_H(t)+y_p(t)
\end{align*}
\begin{itemize}
    \item $y_H(t)$ es la \textcolor{red}{respuesta homogénea} (para $u(t)=0$, con condiciones iniciales)
    \item $y_p(t)$ es la \textcolor{red}{respuesta particular} (para $u(t)\neq0$)
\end{itemize}

\section{Función impulso (Delta Dirac)}
La función \textit{Delta de Dirac} se denota como:
\begin{align*}
    \delta(t)=   \left\{
\begin{array}{ll}
      +\infty & t=0 \\
      0 & t\neq 0 \\
\end{array} 
\right. 
\end{align*}

% Ale: O análogamente.
%\begin{align*}
%    \delta(t-\tau)=   \left\{
%\begin{array}{ll}
%      +\infty & t=\tau \\
%      0 & t\neq \tau \\
%\end{array} 
%\right. 
%\end{align*}

\begin{figure}[h]
\centering
\includegraphics[scale=0.7]{delta.png}
\end{figure}
Propiedades:
\begin{align*}
    \int_{-\infty}^{+\infty}\delta(t)dt=1 \quad &\text{Impulso unitario}\\
    \int_{-\infty}^{+\infty}f(t)\delta(t-T)dt=f(T) \quad &\text{Sampleo}
\end{align*}

\begin{figure}[h]
\centering
\includegraphics[scale=0.7]{dirac.png}
\end{figure}


\begin{align*}
    \Rightarrow f(t)\ast\delta(t-T)=f(T)
\end{align*}

\underline{Propiedad}: Considere que $f(t)$ es una función continua
\begin{align*}
    \int_{-\infty}^{+\infty}f(\tau)\delta(t-\tau)d\tau=f(t) 
\end{align*}
\underline{Importante}: Este resultado permite que cualquier función $f(t)$ pueda igualarse a una suma de impulsos.
\begin{figure}[h]
\centering
\includegraphics[scale=0.7]{suma_imp.png}
\end{figure}

\section{Función de respuesta a impulso unitario}

\begin{figure}[h]
\centering
\includegraphics[scale=0.7]{FRI.png}
\end{figure}

Para sistemas estructurales ($\xi<1$):
\begin{align*}
    m\ddot{x}+c\dot{x}+kx=\delta(t)\\
    x(0)=0 \quad \dot{x}(0)=0
\end{align*}
Solución:
\begin{align*}
    x(t)=h(t)=\frac{1}{m\omega_d}e^{-\xi\omega_nt}\sin{(\omega_dt)}, \quad t\geq0
\end{align*}
Donde $\omega_n = $ y $\omega_d = $. \\

Considere una excitación $u(t)$:
\begin{figure}[h]
\centering
\includegraphics[scale=0.7]{FRI2.png}
\end{figure}
$u(t)$ es continuo $\longrightarrow$ se representa a través de un tren de impulsos $\longrightarrow$ cada impulso tiene una respuesta del sistema $\delta(t-\tau)\longrightarrow h(t-\tau)$.

\begin{figure}[h]
\centering
\includegraphics[scale=0.7]{FRI3.png}
\end{figure}
Por superposición: 
\begin{align*}
    y(t)=\int_0^tu(\tau)h(t-\tau)d\tau \quad \text{``Integral de Duhamel''}\Rightarrow\text{Convolución}
\end{align*}
\underline{En resumen}: Para resolver un sistema LTI, si conoce la $h(t)$, entonces para cualquier $u(t)$, función continua, la respuesta analítica está dada por la convolución.

\begin{align*}
    \Rightarrow y(t)=u(t)\ast h(t)
\end{align*}
\begin{itemize}
    \item $u(t)$: excitación
    \item $\ast$: convolución
    \item $h(t)$: FRI
\end{itemize}

\section{Transformada de Laplace}
\underline{Definición}: La transformada de Laplace de $f(t)$ es igual a:
\begin{align*}
    \mathcal{L}: ~ F(s)=\int_{-\infty}^{+\infty}f(t)e^{(-\xi t)}dt \quad \text{Transformada Bilateral (``two-sided'')}
\end{align*}

donde $s\in\mathbb{C}$, $s=\sigma+j\omega$, $j^2=-1$.\\

\underline{Definición}: La transformada inversa de Laplace es igual a:
\begin{align*}
    \mathcal{L}^{-1}: ~ f(t)=\frac{1}{2\pi j}\int_{+\sigma-j\infty}^{+\sigma+j\infty}F(s)e^{st}ds
\end{align*}
¿Por qué utilizar $\mathcal{L}$? $\Rightarrow$ facilita el análisis de sistemas dinámicos.
\begin{figure}[h]
\centering
\includegraphics[scale=0.7]{laplace.png}
\end{figure}

\underline{Nota}: $H(s)=\mathcal{L}\{h(t) \}$, donde:
\begin{itemize}
    \item $h(t)$: Función de Respuesta Impulso (FRI)
    \item $H(s)$: Función de transferencia (TF)
\end{itemize}

\underline{Aparte}: 
\begin{align*}
    H(j\omega)=\mathcal{F}\{h(t)\}
\end{align*}
\begin{itemize}
    \item $H(j\omega)$: Función de Respuesta en Frecuencia (FRF)
\end{itemize}

\begin{align*}
    H(j\omega) \xleftrightarrow{s=j\omega} H(s) 
\end{align*}

En resumen:
\begin{align*}
    y(t)=h(t)\ast u(t) \xRightarrow[]{\mathcal{L}} Y(s)=H(s)\cdot U(s)
\end{align*}

\begin{align*}
    \mathcal{L}\{y(t)\}=Y(s)=\int_0^{\infty}y(t)e^{-\xi t}dt
\end{align*}

\underline{Propiedades}: 
\begin{itemize}
    \item Superposición: 
    \begin{align*}
        \mathcal{L}\{\alpha f_1(t)+\beta f_1(t)\}=\alpha\mathcal{L}\{f_1(t)\}+\beta\mathcal{L}\{f_2(t)\}
    \end{align*}
    \item Retraso (``time-delay''): 
    \begin{align*}
        f_1(t)&\xrightarrow[]{\mathcal{L}}F_1(s)\\
        f_1(t-\tau)&\xrightarrow[]{\mathcal{L}}e^{-\xi t}F_1(s)
    \end{align*}
    \item ``Time-scaling'': 
    \begin{align*}
        f(\alpha t)&\xrightarrow[]{\mathcal{L}}\frac{1}{\alpha}F(s/\alpha)
    \end{align*}
    \item Diferenciación:
    \begin{align*}
        f(t)&\xrightarrow[]{\mathcal{L}}F(s)\\
        \frac{df(t)}{dt}&\xrightarrow[]{\mathcal{L}}s^2F(s)-sf(0^+)\\
        \frac{d^2f(t)}{dt^2}&\xrightarrow[]{\mathcal{L}}s^2F(s)-sf(0^+)-\dot{f}(0^+)
    \end{align*}
    \item Integración:
    \begin{align*}
        f(t)&\xrightarrow[]{\mathcal{L}}F(s)\\
        g(t)&\xrightarrow[]{\mathcal{L}}G(s)\\
        f(t)\ast g(t)&\xrightarrow[]{\mathcal{L}}F(s)G(s)
    \end{align*}
\end{itemize}




\section{Función de transferencia}
\begin{figure}[h]
\centering
\includegraphics[scale=0.7]{TF.png}
\end{figure}

\begin{itemize}
    \item $H(s)$: Función de transferencia
\end{itemize}

\underline{Definición}: $H(s)\in \mathbb{C}$
\begin{align*}
    H(s)=\frac{\prod_{i=1}^m(s-z_i)}{\prod_{j=1}^n(s-p_j)}
\end{align*}
\begin{itemize}
    \item $z_i$: polos
    \item $p_j$: ceros
\end{itemize}

\underline{Definición}: Polos ($p_j$) son las raíces de la ecuación característica del sistema dinámico:
\begin{align*}
    \prod_{j=1}^n(s-p_j)=0
\end{align*}

\underline{Definición}: Ceros ($z_i$) son las raíces del polinomio del numerador:
\begin{align*}
    \prod_{i=1}^m(s-z_i)=0
\end{align*}

\underline{Importante}: 
\begin{align*}
    H(s)\in \mathbb{C}, \quad s\in\mathbb{C}\\
    \Longrightarrow p_i\in\mathbb{C}\\
    \Longrightarrow z_j \in \mathbb{C}
\end{align*}

Notar que $H(s)$ corresponde a: $H(s)=\mathcal{L}\{h(t)\}$ ($h(t)$ es función de respuesta impulso FRI), por ende, $h(t)=\mathcal{L}^{-1}\{H(s)\}$.\\

\section{Plano complejo}
\begin{itemize}
    \item La función de transferencia recopila \textcolor{red}{toda} la información del sistema dinámico.
    \item Para visualizar las propiedades dinámicas del sistema, vamos a trabajar en el plano complejo.
\end{itemize}

\begin{figure}[h]
\centering
\includegraphics[scale=0.7]{plano_complejo.png}
\end{figure}

$|H(s)|$ representa la ``superficie sobre el plano complejo''
\begin{figure}[h]
\centering
\includegraphics[scale=0.7]{polos_ceros.png}
\end{figure}

\begin{figure}[h]
\centering
\includegraphics[scale=0.7]{magHs.png}
\end{figure}
% Recordar que $H(s)$ corresponde a: $H(s)=\mathcal{L}\{h(t)\}$ ($h(t)$ es FRI),
% Por ende, $h(t)=\mathcal{L}^{-1}\{H(s)\}$.\\

Ejemplo: $H(s)=\frac{1}{s-a}\Longrightarrow h(t)=\mathcal{L}^{-1}\{H(s)\}=e^{at}$\\
\par Por lo tanto, los polos nos dan informaci;on del FRI.\\

\underline{Casos}:
\begin{itemize}
    \item[(i)] $a>0, ~a\in\mathbb{R}\Longrightarrow h(t)=e^{-|a|t}$
    \begin{figure}[h]
    \centering
    \includegraphics[scale=0.7]{caso1.png}
    \end{figure}
    \item[(ii)] $a=0, ~a\in\mathbb{R}\Longrightarrow h(t)=1$
    \begin{figure}[h]
    \centering
    
\includegraphics[scale=0.7]{caso2.png}
    \end{figure}
    \item[(iii)] $a<0, ~a\in\mathbb{R}\Longrightarrow h(t)=e^{|a|t}$
    \begin{figure}[h]
    \centering
    \includegraphics[scale=0.7]{caso3.png}
    \end{figure}
    \item[(iv)] $H(s)=\frac{1}{s^2+a^2}\Longrightarrow s_{1,2}=\pm \beta j \in \mathbb{C}$\\ \\
    $h(t)=\mathcal{L}^{-1}\{H(s)\}=\sin{(\beta t)}$\begin{figure}[h]
    \centering
    \includegraphics[scale=0.7]{caso4.png}
    \end{figure}
    
    \item[(v)] $a_{1,2}=\alpha\pm\beta j \in \mathbb{C}$ (polos conjugados)\\ \\
    \begin{align*}
        H(s) &= \frac{1}{(s-\alpha)^2+\beta^2} \longrightarrow \text{Raíces del denominador}\Rightarrow a_{1,2}\\
        h(t) &= \mathcal{L}^{-1}\{H(s)\} = e^{+\alpha t}\sin{(\beta t)}
    \end{align*}
    \begin{itemize}
        \item[(a)] $\alpha<0$ \\ \\
        $h(t) = e^{-|\alpha| t}\sin{(\beta t)}$
        \begin{figure}[h]
        \centering
        \includegraphics[scale=0.7]{caso5a.png}
        \end{figure}
        
        \item[(b)] $\alpha>0$\\ \\ 
        $h(t) = e^{+|\alpha| t}\sin{(\beta t)}$
        \begin{figure}[h]
        \centering
        \includegraphics[scale=0.7]{caso5b.png}
        \end{figure}
        
        \item[(c)] $\alpha=0$\\ \\
        $h(t)=\sin{(\beta t)}\Longrightarrow$ volvemos al caso (iv).
    \end{itemize}
\end{itemize}

Volviendo al plano complejo, podemos identificar el tipo de respuesta FRI del sistema solo con información de la posición de los polos:
\begin{figure}[h]
\centering
\includegraphics[scale=0.7]{Estabilidad1.png}
\end{figure}

\section{Estabilidad ``Bounded Input, Bounded Output'' (BIBO)}
\begin{figure}[h]
\centering
\includegraphics[scale=0.7]{BIBO.png}
\end{figure}
Un sistema es estable en el sentido BIBO, si y sólo si el sistema posee todos sus polos a la izquierda del plano complejo. Dicho de otra forma: $p_i=\alpha_i\pm\beta_i j$
\begin{itemize}
    \item \textcolor{ForestGreen}{El sistema es estable BIBO ssi $Re(p_i)=\alpha_i<0$}
    \item \textcolor{red}{Si el sistema posee al menos un polo cuyo $Re(p_i)=\alpha_i>0$, entonces el sistema es inestable}
\end{itemize}

\section{Función de transferencia de un sistema de 1GDL}
\begin{figure}[h]
\centering
\includegraphics[scale=0.7]{edm.png}
\end{figure}
\underline{Problema}: Dado EDM, determinar $H(s)$\\ 
\par \underline{Definición}: 
\begin{align*}
    U(s) &= \mathcal{L}\{u(t)\}~: \text{respuesta (dominio complejo)}\\
    F(s) &= \mathcal{L}\{f(t)\} ~: \text{input (dominio complejo)}
\end{align*}

Asumir reposo: $u(0)=\dot{u}(0)=0$

\begin{align*}
    m\ddot{u}+c\dot{u}+ku &= f(t) ~/\color{red}{\mathcal{L}\{\cdot\}}\\
    m\mathcal{L}\{\ddot{u}\}+c\mathcal{L}\{\dot{u}\}+k\mathcal{L}\{u\} &= \mathcal{L}\{f\}\\
    ms^2U(s)+csU(s)+kU(s) &= F(s)\\
    (ms^2+cs+k)U(s) &= F(s)
\end{align*}
\begin{align*}
    \Longrightarrow H(s) = \frac{U(s)}{F(s)}&=\frac{1}{ms^2+cs+k} ~\text{: TF del sistema de 1GDL}
\end{align*}

Aparte: Dado $F(s)$ y $H(s)$
\begin{figure}[h]
\centering
\includegraphics[scale=0.7]{hs.png}
\end{figure}
\begin{align*}
    H(s) &= \frac{1}{ms^2+cs+k} ~\color{red}{\frac{1/m}{1/m}}\\
    &= \frac{1/m}{s^2+2\xi\omega_ns+\omega_n^2}
\end{align*}
\begin{itemize}
    \item $\xi$: razón de amortiguamiento
    \item $\omega_n$: frecuencia natural
\end{itemize}

¿Cómo se obtienen los polos? $\Longrightarrow$ Calcular las raíces del polinomio del denominador:
\begin{align*}
    s^2+2\xi\omega_ns+\omega_n^2=0 ~\text{: Ecuación característica}
\end{align*}
\begin{align*}
     s_{1,2}=-\xi\omega_n\pm\omega_n\sqrt{1-\xi^2}j\\
     s_{1,2}=-\xi\omega_n\pm\omega_dj ~\text{: Polos del sistema de 1GDL}
\end{align*}
¿Es estable? $\Longrightarrow$ Depende del valor de $\xi$:
\begin{itemize}
    \item $\xi>0, ~ Re(s_{1,2}<0\Longrightarrow$ Estable asintóticamente
    \item $\xi<0, ~ Re(s_{1,2}>0\Longrightarrow$ Inestable
    \item $\xi=0, ~ Re(s_{1,2}=0\Longrightarrow$ ``Estable''
\end{itemize}
Además, el sistema oscila sí y solo si $\omega_d>0\Longrightarrow \omega_n>0$: propiedad del sistema dinámico.\\
Si $\xi>0, ~ \omega_n>0$:
\begin{figure}[h]
\centering
\includegraphics[scale=0.6]{estabilidad.png}
\end{figure}

\section{Sistema MGDL}
\begin{align*}
    \tend{U}(s) &= \mathcal{L}\{\tend{u}(t)\}~: \text{respuesta (dominio complejo)}\\
    \tend{F}(s) &= \mathcal{L}\{\tend{f}(t)\} ~: \text{input (dominio complejo)}
\end{align*}
Desde la ecuación de movimiento:
\begin{align*}
    \tenq{M}\tend{\ddot{u}}+\tenq{C}\tend{\dot{u}}+\tenq{K}\tend{u}&=\tend{F} ~/\color{red}{\mathcal{L}\{\cdot\}}\\
    \Longrightarrow\left[s^2\tenq{M}+s\tenq{C}+\tenq{K} \right]\tend{U}(s)&=\tend{F}(s) 
\end{align*}

\underline{Definición}: $\tenq{H}(s)=\left[s^2\tenq{M}+s\tenq{C}+\tenq{K} \right]^{-1}$: TF para un sistema MGDL\\
\par Para buscar los polos, resolver raíces del siguiente polinomio:
\begin{align*}
    s^2\tenq{M}+s\tenq{C}+\tenq{K} = \tend{0} ~/\color{red}{\tend{\cdot \phi}}
\end{align*}
y asumiendo $\tenq{C}=\tenq{0}$
\begin{align*}
    \Longrightarrow (s^2\tenq{M}+\tenq{K})\tend{\phi} = \tend{0} ~\text{: Problema de valores propios generalizado}
\end{align*}



\section{Modelos Dinámicos}
\begin{figure}[h]
\centering
\includegraphics[scale=0.7]{modelos.png}
\end{figure}

La Función de Respuesta en Frecuencia se obtiene a través de:
\begin{align*}
    H(\omega)=\mathcal{F}\{h(t)\}
\end{align*}
De manera alternativa, FRF se obtiene al reemplazar $s=j\omega$ en la TF:
\begin{align*}
    H(\omega)=H(s)\bigg|_{s=j\omega}
\end{align*}
\underline{Ejemplo}: Sistema de 1GDL
\begin{align*}
    \text{EOM: }~m\ddot{u}+c\dot{u}+ku=f(t)
\end{align*}
\begin{align*}
    \Longrightarrow \text{TF: }&~ H(s)=\frac{1}{ms^2+cs+k}=\frac{1/m}{s^2+(c/m)s+k/m}\\
    \Longrightarrow \text{FRF: }&~ H(\omega)=H(s)\bigg|_{s=j\omega}=\frac{1}{(k=m\omega^2)+j\omega c}
\end{align*}
\begin{align*}
    H(\omega)=\frac{1}{(j\omega-\lambda_1)(j\omega-\lambda_1^*)}
\end{align*}
\begin{itemize}
    \item $\{\}^*$: conjugado del número complejo
\end{itemize}
\begin{align*}
    \lambda_1, \lambda_1^*=-\xi\omega_n\pm\omega_n\sqrt{1-\xi^2}~\text{: polos del sistema}
\end{align*}

\subsection{Representación en Variables de Estado}
\begin{align*}
    \text{EDM: }~\tenq{M}\tend{\ddot{u}}+\tenq{C}\tend{\dot{u}}+\tenq{K}\tend{u}&=\tenq{G}\tend{p}(t) 
\end{align*}
\begin{itemize}
    \item $\tenq{G}$: matriz de influencia
\end{itemize}

\underline{Definición}: 
\begin{itemize}
    \item Vector de Estado: $\tend{x}=\begin{bmatrix}\tend{u}\\ \tend{\dot{u}}\end{bmatrix}$
    \item Ecuación de Estado: $\tend{\dot{u}}=\tenq{A_s}\tend{x}+\tenq{B_s}\tend{p}$
\end{itemize}

\begin{align*}
    \Longrightarrow \tend{\dot{x}}=\begin{bmatrix}\tend{\dot{u}}\\ \tend{\ddot{u}}\end{bmatrix}=\tenq{A_s}\begin{bmatrix}\tend{u}\\ \tend{\dot{u}}\end{bmatrix}+\tenq{B_s}\tend{p}(t)
\end{align*}

Donde:
\begin{itemize}
    \item Matriz de estado: $\tenq{A_s}=\begin{bmatrix}\tenq{0} & \tenq{I}\\ -\tenq{M}^{-1}\tenq{K} & -\tenq{M}^{-1}\tenq{C}\end{bmatrix}$
    \item Matriz de influencia de excitación: $\tenq{B_s}=\begin{bmatrix}\tenq{0}\\ \tenq{M}^{-1}\tenq{G}\end{bmatrix}$
\end{itemize}

\underline{Definición}: Ecuación de respuesta\\
\begin{itemize}
    \item Vector de respuesta: $\tend{y}$
\end{itemize}

\begin{align*}
    \tend{y}=\tenq{C_s}\tend{x}+\tenq{D_s}\tend{p}
\end{align*}

\underline{Ejemplo}: $\tend{y}=\begin{bmatrix}\tend{u} \\ \tend{\dot{u}}\\ \tend{\ddot{u}}\end{bmatrix}$

\begin{itemize}
    \item $\tenq{C_s}=\begin{bmatrix}\tenq{I} & \tenq{0}\\
    \tenq{0} & \tenq{I}\\ 
    -\tenq{M}^{-1}\tenq{K} & -\tenq{M}^{-1}\tenq{C}\end{bmatrix}$
    \item ``Feedthrough'': $\tenq{D_s}=\begin{bmatrix}\tenq{0}\\
    \tenq{0}\\ \tenq{M}^{-1}\tenq{G}\end{bmatrix}$
\end{itemize}

\underline{Resumen}: 

\begin{align*}
\left .
\begin{array}{rr}
\tend{\dot{x}}=\tenq{A_s}\tend{x}+\tenq{B_s}\tend{p}\\
\tend{y}=\tenq{C_s}\tend{x}+\tenq{D_s}\tend{p} 
\end{array}
\right \} \Longrightarrow 
\begin{array} {|r|r|}
  \hline A & B \\
  \hline C & D \\
  \hline 
\end{array}
\end{align*}

\begin{figure}[h]
\centering
\includegraphics[scale=0.7]{SS.png}
\end{figure}

\subsection{Relación entre State-Space y TF}
Asumir condiciones iniciales: $\tend{x}(0)=0\Longrightarrow$ Reposo.\\
\par \underline{Definición}: $\mathcal{L}\{\tend{x}(t)\}=\tend{X}(s)$, $\mathcal{L}\{\tend{p}(t)\}=\tend{P}(s)$ \\
Entonces: 
\begin{align*}
    \tend{\dot{x}}&=\tenq{A_s}\tend{x}+\tenq{B_s}\tend{p} ~/\color{red}{\mathcal{L}\{\cdot\}}\\
    \Longrightarrow s\tend{X}(s)&= \tenq{A_s}\tend{X}(s)+\tenq{B_s}\tend{P}(s)\\
    \Longrightarrow (s\tenq{I}-\tenq{A_s})\tend{X}(s)&=\tenq{B_s}\tend{P}(s)
\end{align*}
\begin{align}
    \Longrightarrow \tend{X}(s)&=(s\tenq{I}-\tenq{A_s})^{-1}\tenq{B_s}\tend{P}(s)
\end{align}

\par Ecuación de respuesta:
\par \underline{Definición}: $\mathcal{L}\{\tend{y}(t)\}=\tend{Y}(s)$\\
\begin{align*}
    \tend{y}(t)=\tenq{C_s}\tend{x}+\tenq{D_s}\tend{p} ~/\color{red}{\mathcal{L}\{\cdot\}}\\
\end{align*}
\begin{align}
    \Longrightarrow \tend{Y}(s)&=\tenq{C_s}\tend{X}(s)+\tenq{D_s}\tend{P}(s)
\end{align} 

Reemplazando la ecuación (2) en (1):
\begin{align*}
    \tend{Y}(s)&=\tenq{C_s}(s\tenq{I}-\tenq{A_s})^{-1}\tenq{B_s}\tend{P}(s)+\tenq{D_s}\tend{P}(s)\\
    \tend{Y}(s)&=[\tenq{C_s}(s\tenq{I}-\tenq{A_s})^{-1}\tenq{B_s}+\tenq{D_s}]\tend{P}(s)\\
    \Longrightarrow \tenq{H}(s)&=\tenq{C_s}(s\tenq{I}-\tenq{A_s})^{-1}\tenq{B_s}+\tenq{D_s}
\end{align*}
Asumiendo $\tenq{D_s}=0$:
\begin{align*}
\tenq{H}(s)&=\tenq{C_s}(s\tenq{I}-\tenq{A_s})^{-1}\tenq{B_s}\\
\tenq{H}(s)&=\frac{1}{\det(s\tenq{I}-\tenq{A_s})}\tenq{C_s}\text{ adj}(s\tenq{I}-\tenq{A_s})\tenq{B_s}
\end{align*}

\begin{align*}
    \Longrightarrow\det(s\tenq{I}-\tenq{A_s})=0 ~\text{: Problema de valores propios de }\tenq{A_s}
\end{align*}

Polos:
\begin{itemize}
    \item Raíces del polinomio del denominador de $\tenq{H}(s)$
    \item Valores propios de $\tenq{A_s}$
\end{itemize}

\section{Implementación utilizando Matlab/Simulink}
Matlab $\longrightarrow$ variables del tipo sistema\\
\par $\tenq{M}$, $\tenq{C}$, $\tenq{K}$, $\tenq{G} \Longrightarrow$ matrices de tipo ``double'' \\

\par $\Longrightarrow$ $\tenq{A_s}$, $\tenq{B_s}$, $\tenq{C_s}$, $\tenq{D_s} \Longrightarrow$ matrices de tipo ``double''\\
\par Para construir un sistema de tipo espacio-estado, se utiliza el siguiente comando:
\begin{minted}[]{matlab}
>> sys = ss(As, Bs, Cs, Ds); % Variable tipo sistema 
\end{minted}

\begin{align*}
    H(s)=\frac{b_ms^m+b_{m-1}s^{m-1}+...+b_1s+b_0}{s^n+a_{n-1}s^{n-1}+...+a_1s+a_0}
\end{align*}

\begin{minted}[]{matlab}
>> sys = tf(num,den); % Variable sistema 
\end{minted}
Donde:
\begin{align*}
    \text{num} = [b_m, b_{m-1}, ... , b_1, b_0]\\
    \text{den} = [1, a_{n-1}, ... , a_1, a_0]\\
\end{align*}
\underline{Ejemplo}: 1GDL
\begin{align*}
    \text{num} = [1/m]\\
    \text{den} = [1, c/m, k/m]\\
\end{align*}

\underline{Matlab: Control System Toolbox}\\
\begin{itemize}
    \item Valores y vectores propios del sistema:
\begin{minted}[]{matlab}
>> eig(sys)
>> eig(As)
\end{minted}
\item Polos del sistema, frecuencias naturales y razones de amortiguamiento:
\begin{minted}[]{matlab}
>> damp(sys)
>> damp(As)
\end{minted}
\item Posición de ``ceros'' y ``polos'' en el plano complejo:
\begin{minted}[]{matlab}
>> pzmap(sys)
\end{minted}
\item Gráfico magnitud y fase vs frecuencia del sistema
\begin{minted}[]{matlab}
>> bode(sys)
>> bodeplot(sys)
\end{minted}
\item Definición de un sistema con su TF:
\begin{align*}
    \text{TF: }~H(s)= k\frac{\prod_{i=1}^m(s-z_i)}{\prod_{j=1}^n(s-p_j)}
\end{align*} 
\begin{align*}
    \text{zeros} &= [z_m, z_{m-1}, ..., z_2, z_1]\\
    \text{poles} &= [p_n, p_{n-1},..., p_2, p_1]\\
    \text{gain} &= k
\end{align*} 
\begin{minted}[]{matlab}
>> sys = zpk(zeros, polos, gain); 
\end{minted}

\item Ganancia del sistema
\begin{minted}[]{matlab}
>> dcgain(sys) %% DC: direct current (caso estático)
\end{minted}

\item Análisis: respuesta a impulso y a escalón
\begin{minted}[]{matlab}
>> impulse(sys)
>> step(sys)
\end{minted}
\end{itemize}




\section{Función respuesta escalón (``step'')}
\begin{figure}[h]
\centering
\includegraphics[scale=0.7]{step.png}
\end{figure}
\underline{Definición}:
\begin{itemize}
    \item \textcolor{red}{Tiempo de subida} (``rise time'', $t_r$): tiempo que toma ir desde 10\% hasta 90\% de respuesta (alternativa: 0\% a 100\%)
    \item \textcolor{Goldenrod}{Sobreimpulso} (``overshoot'', $M_p$)
    \item \textcolor{LimeGreen}{Tiempo de asentamiento} (``senttling time'', $t_s$): tiempo que toma el sistema en decaer hasta 1\%
\end{itemize}

\section{Especificaciones de Respuesta Transiente}
Suponiendo un sistema LTI-SISO:\\
\begin{align*}
    \text{EDM: }~ m\ddot{u}+c\dot{u}+ku &= f\\
    \Longrightarrow H(s) &= \frac{1/m}{s^2+2\xi\omega_n s+\omega_n^2}
\end{align*}
\begin{figure}[h]
\centering
\includegraphics[scale=0.5]{transiente.png}
\end{figure}

Fórmulas para estimar la respuesta transiente para el caso de excitación escalón:
\begin{itemize}
    \item \underline{Tiempo subida, $t_r$}: Asumiento $\xi=2\%$, 
    $$t_r\approx \frac{1.6}{\omega_n}$$
    \item \underline{Sobreimpulso, $M_p$}: 
    $$-\log \left(\frac{M_p}{\pi} \right) \approx \tan(\theta)$$
    \item \underline{Tiempo asentamiento, $t_s$}: Asumiendo 2\% de respuesta esperada en ``steady-state'',
    $$t_s\approx \frac{4}{\xi \omega_n} $$
\end{itemize}

\section*{``Pole placement'', ubicación de polos}
Dadas las especificaciones transientes: $t_r$, $M_p$, $t_s$
\begin{figure}[h]
\centering
\includegraphics[scale=0.7]{pole_placement.png}
\end{figure}

\begin{align*}
    \color{VioletRed}{t_s}&<(t_s)_{lim}=\frac{4}{(\xi\omega_n)_{lim}}\\
    \color{BurntOrange}{t_r}&<(t_r)_{lim}\approx \frac{1.6}{\omega_n}\\
    \color{LimeGreen}{M_p}&<(M_p)_{lim}=e^{-\pi\tan(\theta)}
\end{align*}
